% Preâmbulo LaTeX do relatório (entra no cabeçalho YAML do RELATORIO.md, em header-includes).
% Títulos a negrito: 14 pt (nível 1), 12 pt (nível 2), 11 pt (nível 3).
\usepackage{titlesec}
\titleformat*{\section}{\fontsize{14}{17}\selectfont\bfseries}
\titleformat*{\subsection}{\fontsize{12}{14.5}\selectfont\bfseries}
\titleformat*{\subsubsection}{\normalsize\bfseries}
% Legendas por cima, "Tabela n." / "Figura n." a negrito e o texto em itálico.
\usepackage{floatrow}
\floatsetup[figure]{capposition=top}
\floatsetup[table]{capposition=top}
\usepackage{caption}
\captionsetup{labelfont=bf,textfont=it,labelsep=period,justification=justified,singlelinecheck=false}
% Tabelas a 10 pt, cabeçalho cinzento claro, linhas horizontais finas.
\usepackage{etoolbox}
\AtBeginEnvironment{longtable}{\small}
\usepackage{colortbl}
\definecolor{cabecalho}{gray}{0.88}
\setlength{\heavyrulewidth}{0.6pt}
\setlength{\lightrulewidth}{0.4pt}
\let\toprulebase\toprule
\renewcommand{\toprule}{\toprulebase\rowcolor{cabecalho}}
% Índice até ao nível 2; espaço simples depois do ponto final.
\setcounter{tocdepth}{2}
\frenchspacing
